\chapter{多重积分}
\label{chap:multiple-integrals-solutions}

\begin{solution}{multiple-integral-product-function}
设 $|f|\le M$、$|g|\le N$。若某区间退化，两端均为零。以下设
$a<b,c<d$，先证明 $h(x,y)=f(x)g(y)$ 可积。对乘积分割
$P=P_x\times P_y$，在 $I_i\times J_j$ 上有
\[
 \omega(h;I_i\times J_j)
 \le N\omega(f;I_i)+M\omega(g;J_j).
\]
乘以 $|I_i||J_j|$ 并求和，得到
\[
 U(h,P)-L(h,P)
 \le N(d-c)\sum_i\omega(f;I_i)|I_i|
     +M(b-a)\sum_j\omega(g;J_j)|J_j|.
\]
分别选择一元振幅和足够小的分割，右端可任意小，所以 $h$ 可积。
这一步已保证任意标记的细分割和都趋向同一积分。现在为求这个值，
特选可分离标记 $(\xi_i,\eta_j)$，便有
\[
 \sum_{i,j}f(\xi_i)g(\eta_j)\Delta x_i\Delta y_j
 =\left(\sum_i f(\xi_i)\Delta x_i\right)
  \left(\sum_j g(\eta_j)\Delta y_j\right).
\]
令两个分割的网径同时趋于零，右端两个因子分别趋于一元积分，故
\[
 \int_{[a,b]\times[c,d]}f(x)g(y)\,d(x,y)
 =\int_a^bf(x)\,dx\int_c^dg(y)\,dy.
\]
仅有可分离标记下的等式不足以证明可积性；前面的振幅估计正是用来处理
各块标记彼此独立的情形。
\end{solution}

\begin{solution}{multiple-continuous-integrability}
设 $Q\subset\R^n$ 是非退化闭矩形。连续函数在紧集 $Q$ 上一致连续，故给定
$\varepsilon>0$，存在 $\delta>0$，使同一子矩形中任意两点距离小于
$\delta$ 时函数值之差小于 $\varepsilon/(2|Q|)$。取网径充分小的分割，
每个子矩形 $Q_j$ 上的振幅 $\omega_j$ 至多为该数，因而
\[
 U(f,P)-L(f,P)=\sum_j\omega_j|Q_j|\le\varepsilon/2<\varepsilon.
\]
Darboux 判据给出可积性。若 $Q$ 有一个边长为零，则每个子矩形的 $n$ 维
体积均为零，上和与下和都为零；按照本章约定积分为零。
\end{solution}

\begin{solution}{jordan-boundary-indicator}
若 $Q$ 退化，则 $\partial E\subseteq Q$ 是零测集，指示函数积分为零，
结论成立。以下设 $Q$ 非退化。若 $\partial E$ Jordan 零测，则指示函数
在 $Q$ 上的相对不连续点包含于 $\partial E$，由零测不连续集判据可积。

反向设指示函数可积。给定 $\varepsilon>0$，取分割 $P$，使
\[
 \sum_{R\in P}\omega(\mathbf1_E;R)|R|<\varepsilon/2.
\]
令 $B$ 为所有分割面的并（包含 $\partial Q$），以总体积小于
$\varepsilon/2$ 的有限薄矩形覆盖 $B$。任何 $x\in\partial E\setminus B$
都处在某一块内部；在该块内的小球同时遇到 $E$ 与补集，所以该块的振幅
为一。这些振幅为一的块，加上覆盖 $B$ 的薄矩形，覆盖整个 $\partial E$，
总体积小于 $\varepsilon$。因此 $\partial E$ Jordan 零测。

必须保留 $B$：例如 $E=Q$ 时，指示函数在 $Q$ 上处处连续，却有
$\partial E=\partial Q\ne\varnothing$。两个外包矩形可嵌入同一个大矩形，
把各矩形面加入分割，零延拓与矩形分割可加性保证积分值不变。
\end{solution}

\begin{solution}{graph-jordan-null}
设 $f:[a,b]\to\R$ 的 Lipschitz 常数为 $L$。把 $[a,b]$ 等分为 $N$ 段，
长度 $h=(b-a)/N$。每段上的函数振幅不超过 $Lh$，故该段上方的图像可由
一个宽 $h$、高 $2Lh$（端点稍加任意小余量）的矩形覆盖。总面积不超过
$2LN h^2=2L(b-a)^2/N$，趋于零。若 $L=0$，同样可用任意小高度的矩形条
覆盖。故图像为 Jordan 零测集。
\end{solution}

\begin{solution}{dense-countable-not-jordan-null}
若有限个闭矩形覆盖 $\Q^n\cap[0,1]^n$，其并集是闭集；由于有理点在立方体
中稠密，该并集也覆盖其闭包 $[0,1]^n$。所以这些矩形体积之和至少为 $1$，
不可能任意小，故该集合不是 Jordan 零测集。Lebesgue 零测允许可数个区间
覆盖：逐个有理点用体积总和小于 $\varepsilon$ 的小立方体覆盖即可。Jordan
理论只允许有限覆盖，两种结论因覆盖族不同而不冲突。
\end{solution}

\begin{solution}{jordan-volume-finite-union}
先证两个集合。指标恒等式
\[
 \mathbf1_{E\cup F}=\mathbf1_E+\mathbf1_F-\mathbf1_{E\cap F}
\]
处处成立；各指标函数均可积，因为
$\partial(E\cap F),\partial(E\cup F)\subseteq\partial E\cup\partial F$。
积分的线性性给出
\[
 |E\cup F|=|E|+|F|-|E\cap F|.
\]
若 $E^\circ\cap F^\circ=\varnothing$，而 $x\in E\cap F$，则 $x$ 不可能
同时是 $E$ 与 $F$ 的内点。若它是 $E$ 的内点，则它不是 $F$ 的内点，且
$x\in F$，所以 $x\in\partial F$；另一种情形同理。因此
\[
 E\cap F\subseteq\partial E\cup\partial F.
\]
右端为 Jordan 零测集，故 $|E\cap F|=0$，从而
$|E\cup F|=|E|+|F|$。对于有限族 $(E_i)_{i=1}^m$，每个两两交集
$E_i\cap E_j$ 都是 Jordan 零测集，因而
\[
 E_k\cap\bigcup_{i<k}E_i=\bigcup_{i<k}(E_k\cap E_i)
\]
也是 Jordan 零测集。逐次使用两个集合的公式即得
$|\bigcup_{i=1}^mE_i|=\sum_{i=1}^m|E_i|$。
\end{solution}

\begin{solution}{multiple-darboux-cauchy}
设 $|f|\le M$，先设 $Q$ 非退化且 $M>0$。若振幅判据成立，给定
$\varepsilon>0$，选 $P_0$ 使 $U(f,P_0)-L(f,P_0)<\varepsilon/3$。
记其有限个内部分割面附近、宽 $2\eta$ 的薄条总体积上界为 $C\eta$。
选 $\eta$ 使 $4MC\eta<\varepsilon/3$。对任意网格小于 $\eta$ 的
分割 $P$，将其加细为同时包含 $P_0$ 的分割。没有碰到 $P_0$ 分割面的
块可以保留原标记；被切开的块总体积至多为 $C\eta$。在其中任取新标记，
新旧和之差至多为 $2MC\eta$。加细和介于 $L(f,P_0)$ 与 $U(f,P_0)$
之间。因此任意两个网格小于 $\eta$ 的标记和相差至多
\[
 U(f,P_0)-L(f,P_0)+4MC\eta<\varepsilon.
\]
这就是 Cauchy 判据，其中两个分割及其标记均可独立选择。

反之，给定 $\varepsilon>0$，按 Cauchy 条件选网格阈值，使任意两个
充分细标记和相差小于 $\varepsilon/2$。固定这样的分割 $P$，分别在每块
选取逼近上、下确界的标记；让逼近误差趋零，可得
$U(f,P)-L(f,P)\le\varepsilon/2<\varepsilon$。于是振幅判据成立。

为同时核验 Darboux 判据与 Riemann 定义的关系，公共加细使每个下和
不超过每个上和。振幅和可任意小时，上下积分相等；上面的薄条估计说明
所有充分细的标记和均趋于公共值。若反过来所有细标记和趋于同一值，则
它们满足 Cauchy 条件，前段给出振幅判据。$M=0$ 或 $Q$ 退化时所有和
均为零，结论直接成立。
\end{solution}

\begin{solution}{jordan-set-algebra}
对任意集合 \(E,F\)，有
\[
 \partial(E\cup F),\ \partial(E\cap F),\ \partial(E\setminus F)
 \subseteq\partial E\cup\partial F.
\]
例如，若一点既不在 \(\partial E\) 也不在 \(\partial F\)，则它有一个
邻域，在其中对 \(E\)、\(F\) 的成员关系分别恒定；并、交、差的成员关系也
恒定，所以该点不在相应边界。若 \(E,F\) Jordan 可测，则右端是有限个
Jordan 零测集之并，仍为 Jordan 零测集。集合判据于是说明三种集合运算的
结果都 Jordan 可测。
\end{solution}

\begin{solution}{jordan-product-volume}
集合恒等式
\[
 \partial(E\times F)
 \subseteq(\partial E\times\overline F)
       \cup(\overline E\times\partial F)
\]
成立。把 \(\partial E\)、\(\partial F\) 分别用总体积任意小的有限矩形
覆盖，再用包住 \(E,F\) 的固定矩形控制另一因子，便知右端是
\(\R^{p+q}\) 中的 Jordan 零测集。因此 \(E\times F\) Jordan 可测。

指标函数分解为
\(\mathbf1_{E\times F}(x,y)=\mathbf1_E(x)\mathbf1_F(y)\)。对乘积分割，
相应 Riemann 和分解成两个因子。令网径趋于零得到
\[
 |E\times F|=|E|\,|F|.
\]
\end{solution}

\begin{solution}{jordan-continuous-domain}
$\overline E$ 紧，所以 $f$ 有界。对包含 $E$ 的闭矩形 $Q$，令
$f_E=f$ 于 $E$，并令 $f_E=0$ 于 $Q\setminus E$。若一点不在
$\partial E$ 上，则它或者在 $E$ 的内部，或者在 $\overline E$ 的外部；
前一种情形由 $f$ 的连续性，后一种情形由局部恒零性，说明 $f_E$ 连续。
因此不连续点包含于 Jordan 零测集 $\partial E$，从而 $f_E$ 可积。

若 $F$ 是 $f|_E$ 的任意延拓，不要求它连续，则仍有 $F\mathbf1_E=f_E$，
所以积分值与延拓选择无关。对另一个外包矩形 $Q'$，取同时包含 $Q,Q'$
的大矩形并按二者的面分割。零延拓在 $Q$ 外为零，矩形分割可加性给出
大矩形积分等于 $Q$ 上的积分；对 $Q'$ 同理。这证明外包矩形无关性。
\end{solution}
